\section{Dyadic irreducibility and exact capacity limitations}
\label{sec:dyadic-capacity}

\subsection{Dyadic irreducibility and minimal Weil-height saturation}
\label{subsec:dyadic-irreducibility}

The defining-height result above leaves open, for a general truncation
index, the possibility that the root selected near a zero of \(\Psi\)
belongs to a factor of much smaller degree or height.  On the dyadic
subsequence this possibility can be excluded completely.  Recall
\(H_N\), \(D_N^*\), and \(Q_N\) from
\eqref{eq:HN-truncation} and \eqref{eq:QN-definition}, and put
\begin{equation}\label{eq:PN-AN-dyadic}
 P_N(w)=D_N^*H_N(w),
 \qquad
 A_N=\frac{D_N^*}{N\,N!}.
\end{equation}
Thus, with \(b_{m,N}=D_N^*/(m\,m!)\),
\begin{align}
 P_N(w)
 &=\sum_{m=1}^N(-1)^{m+1}b_{m,N}C_m(w),
 \label{eq:PN-chebyshev}\\
 Q_N(z)
 &=z^NP_N(z+z^{-1})
   =\sum_{m=1}^N(-1)^{m+1}b_{m,N}
       \bigl(z^{N+m}+z^{N-m}\bigr).
 \label{eq:QN-dyadic-explicit}
\end{align}

\begin{theorem}[Dyadic Eisenstein and reciprocal irreducibility]
\label{thm:dyadic-reciprocal-irreducibility}
Let \(N=2^r\geq2\).  Then \(A_N\) is an odd integer and
\begin{equation}\label{eq:PN-mod-four}
 \boxed{P_N(w)\equiv-A_N(w^N+2)\pmod 4.}
\end{equation}
Consequently \(P_N\) is a primitive \(2\)-Eisenstein polynomial and is
irreducible over \(\Q\) of degree \(N\).  Moreover, \(Q_N\) is primitive
and irreducible over \(\Q\) of degree \(2N\).
\end{theorem}

\begin{proof}
Legendre's formula gives
\[
 v_2(N!)=N-s_2(N)=N-1,
\]
where \(s_2(N)\) is the sum of the binary digits of \(N\).  Hence
\begin{equation}\label{eq:dyadic-top-valuation}
 v_2(N\,N!)=N+r-1.
\end{equation}
If \(1\leq m<N\), then \(v_2(m)\leq r-1\) and
\(v_2(m!)\leq m-1\), so
\[
 v_2(m\,m!)\leq m+r-2\leq N+r-3.
\]
It follows that the maximum defining \(v_2(D_N^*)\) is attained at
\(m=N\), that \(A_N\) is odd, and that
\begin{equation}\label{eq:lower-dyadic-coefficients-four}
 4\mid b_{m,N}\qquad(1\leq m<N).
\end{equation}

The identity \(C_{2m}=C_m^2-2\), beginning with
\(C_2(w)=w^2-2\), gives by induction
\begin{equation}\label{eq:dyadic-chebyshev-mod-four}
 C_{2^r}(w)\equiv w^{2^r}+2\pmod4.
\end{equation}
Since \(N\) is even, the coefficient of \(C_N\) in
\eqref{eq:PN-chebyshev} is \(-A_N\).  Equations
\eqref{eq:lower-dyadic-coefficients-four} and
\eqref{eq:dyadic-chebyshev-mod-four} prove
\eqref{eq:PN-mod-four}.  Thus the leading coefficient of \(P_N\) is odd,
all its remaining coefficients are even, and its constant coefficient is
congruent to \(-2A_N\) modulo \(4\).  The polynomial \(Q_N\) is primitive
by Lemma~\ref{lem:truncation-denominator-scale}; hence \(P_N\) is
primitive as well, because any common divisor of the coefficients of
\(P_N\) would divide every coefficient of \(z^NP_N(z+z^{-1})=Q_N\).
Eisenstein's criterion now proves the irreducibility of \(P_N\).

It remains to check that the reciprocal substitution does not split the
polynomial.  Since \(C_m(2)=2\) and \(C_m(-2)=2(-1)^m\),
\begin{align}
 P_N(2)
 &=2D_N^*\sum_{m=1}^N\frac{(-1)^{m+1}}{m\,m!}>0,
 \label{eq:PN-at-two-positive}\\
 P_N(-2)
 &=-2D_N^*\sum_{m=1}^N\frac{1}{m\,m!}<0.
 \label{eq:PN-at-minus-two-negative}
\end{align}
The first sign follows by pairing consecutive terms of the decreasing
alternating sum, since \(N\) is even.

Let \(\theta\) be a root of \(P_N\), let
\(K=\Q(\theta)\), and write \(a_N=\operatorname{lc}(P_N)=-A_N\).
Because \(P_N\) is irreducible of degree \(N\),
\begin{equation}\label{eq:dyadic-discriminant-norm}
 \operatorname N_{K/\Q}(\theta^2-4)
 =\frac{P_N(2)P_N(-2)}{a_N^2}<0.
\end{equation}
A square in \(K\) has norm equal to a square in \(\Q\), so
\(\theta^2-4\) is not a square in \(K\).  Therefore
\(X^2-\theta X+1\) is irreducible over \(K\).  If \(\alpha\) is one of
its roots, then \(\theta=\alpha+\alpha^{-1}\) and consequently
\[
 [\Q(\alpha):\Q]
 =2[K:\Q]=2N.
\]
The number \(\alpha\) is a zero of \(Q_N\), whose degree is \(2N\).
Thus \(Q_N\) is irreducible; its primitivity was already established in
Lemma~\ref{lem:truncation-denominator-scale}.
\end{proof}

\begin{theorem}[Dyadic minimal-height saturation]
\label{thm:dyadic-minimal-height-saturation}
Let \(\rho\) be a fixed simple zero of \(\Psi\) with \(|\rho|>1\), and,
for all sufficiently large \(N\), let \(\alpha_N\) be the zero of
\(\Psi_N\) converging to \(\rho\) in
Theorem~\ref{thm:truncation-height-saturation}.  Along
\(N=2^r\to\infty\),
\begin{equation}\label{eq:dyadic-degree-height}
 [\Q(\alpha_N):\Q]=2N,
 \qquad
 h(\alpha_N)=\frac12\log N+O(1).
\end{equation}
More precisely,
\begin{align}
 \frac{-\log|\alpha_N-\rho|}
 {[\Q(\alpha_N):\Q]h(\alpha_N)}&\longrightarrow1,
 \label{eq:dyadic-root-height-ratio}\\
 \frac{-\log|\Psi(\alpha_N)|}
 {[\Q(\alpha_N):\Q]h(\alpha_N)}&\longrightarrow1.
 \label{eq:dyadic-value-height-ratio}
\end{align}
Here \(h\) is the absolute logarithmic Weil height.  In particular,
\(Q_N\), up to sign, is the primitive integer minimal polynomial of
\(\alpha_N\); the full clearing polynomial is not concealing a factor of
smaller degree or height on the dyadic subsequence.
\end{theorem}

\begin{proof}
By Theorem~\ref{thm:dyadic-reciprocal-irreducibility}, \(Q_N\) is the
primitive irreducible integer polynomial of \(\alpha_N\), up to sign.
The coefficient--Mahler inequalities for a polynomial of degree \(2N\),
together with Lemma~\ref{lem:truncation-denominator-scale}, give
\begin{equation}\label{eq:dyadic-mahler-bounds}
 4^{-N}D_N^*
 \leq M(Q_N)
 \leq \lVert Q_N\rVert_2
 \leq \sqrt{2N+1}\,D_N^*.
\end{equation}
Indeed, the lower bound follows from
\(H(Q_N)\leq2^{2N}M(Q_N)\), while the upper bound is Landau's inequality
and \(H(Q_N)=D_N^*\).  Since
\(\log D_N^*=N\log N+o(N)\), it follows that
\begin{equation}\label{eq:dyadic-mahler-asymptotic}
 \log M(Q_N)=N\log N+O(N).
\end{equation}
The degree assertion in \eqref{eq:dyadic-degree-height} follows from
Theorem~\ref{thm:dyadic-reciprocal-irreducibility}, and the standard height
identity gives
\[
 2N\,h(\alpha_N)=\log M(Q_N).
\]
This proves the height estimate in \eqref{eq:dyadic-degree-height}.

The root asymptotic \eqref{eq:truncation-root-asymptotic} gives
\begin{equation}\label{eq:dyadic-root-log-scale}
 -\log|\alpha_N-\rho|=N\log N+O_\rho(N).
\end{equation}
Moreover \(\Psi_N(\alpha_N)=0\), so the uniform tail estimate used in
the proof of Theorem~\ref{thm:truncation-height-saturation} yields
\[
 \Psi(\alpha_N)
 =\frac{(-1)^{N+2}\alpha_N^{N+1}}
 {(N+1)(N+1)!}
 \left(1+O_\rho(N^{-1})\right).
\]
Since \(\alpha_N\to\rho\),
\begin{equation}\label{eq:dyadic-value-log-scale}
 -\log|\Psi(\alpha_N)|=N\log N+O_\rho(N).
\end{equation}
Dividing \eqref{eq:dyadic-root-log-scale} and
\eqref{eq:dyadic-value-log-scale} by
\(2N h(\alpha_N)=N\log N+O(N)\) proves
\eqref{eq:dyadic-root-height-ratio} and
\eqref{eq:dyadic-value-height-ratio}.
\end{proof}

\begin{corollary}[The precise moving-point threshold]
\label{cor:dyadic-moving-point-threshold}
For every real \(c<1\), infinitely many algebraic numbers in the dyadic
family above satisfy
\begin{equation}\label{eq:dyadic-lower-bound-failure}
 \log|\Psi(\alpha_N)|
 <-c\,[\Q(\alpha_N):\Q]h(\alpha_N).
\end{equation}
Equivalently, no estimate of the form
\[
 -\log|\Psi(\alpha)|
 \leq(c+o(1))[\Q(\alpha):\Q]h(\alpha),
 \qquad c<1,
\]
can hold uniformly on any moving family containing these dyadic points.
For every \(c>1\), the reverse strict inequality to
\eqref{eq:dyadic-lower-bound-failure} holds eventually along this family,
so this sequence alone cannot contradict a bound with coefficient greater
than one.
\end{corollary}

\begin{proof}
This is exactly the limit in
\eqref{eq:dyadic-value-height-ratio}.
\end{proof}

\begin{remark}[Scope of the obstruction]
The corollary identifies the normalized threshold for the term
\([\Q(\alpha):\Q]h(\alpha)\) alone.  It does not rule out estimates with
additional terms of comparable size, such as
\([\Q(\alpha):\Q]\log[\Q(\alpha):\Q]\), nor does it exclude nonlinear
auxiliary constructions or non-dyadic phenomena.  What it closes is the
specific escape through a smaller irreducible factor or smaller minimal
Weil height for the selected truncation root.
\end{remark}
